
@@M0@@

\textbf{Remarque 1.14.} En utilisant la notion de «schéma relatif» (cf [3]), on peut encore donner un sens à 1.1 en supposant que @@M1@@ est un topos localement annelé. Si on suppose que l'objet final du topos est quasi-compact, la démonstration donnée dans le cas des schémas se transpose sans difficultés, et le théorème reste donc valable.

\section*{2. Calcul du @@M2@@ d'un fibré projectif : cas des complexes parfaits}


Le présent paragraphe ne sera pas utilisé dans la suite du séminaire.

La notation @@M3@@ désignera dans ce paragraphe le groupe de Grothendieck des complexes parfaits sur @@M4@@ (IV 2.2). Nous nous proposons de prouver le théorème suivant, analogue de 1.1 :

\textbf{Théorème 2.1.} \uline{Soient} @@M5@@ \uline{un schéma quasi-compact}, @@M6@@ \uline{un} @@M7@@\uline{-Module localement libre de rang} @@M8@@, @@M9@@ \uline{le fibré projectif associé} ; \uline{soit} @@M10@@ \uline{la classe} \uline{dans} @@M11@@ \uline{du complexe égal à} @@M12@@ \uline{en degré nul}, \uline{et nul pour les autres degrés}. \uline{Alors} @@M13@@, \uline{muni de sa structure naturelle de} @@M14@@-\uline{algèbre est} \uline{un} @@M15@@-\uline{module libre admettant pour base} @@M16@@.

2.2. Fixons les notations : @@M17@@ désignera l'algèbre symétrique de @@M18@@; on notera @@M19@@ la catégorie des @@M20@@-Modules, @@M21@@ la catégorie des @@M22@@-Modules gradués avec pour morphismes les morphismes gradués, @@M23@@ la sous-catégorie pleine de @@M24@@ dont les objets sont les @@M25@@-Modules localement libres de type fini en tant que @@M26@@-Modules gradués, @@M27@@ la catégorie des @@M28@@-Modules. Les catégories @@M29@@, @@M30@@, et @@M31@@ sont abéliennes.

Pour tout @@M32@@, et tout complexe @@M33@@ de @@M34@@, nous noterons @@M35@@ le complexe @@M36@@.

